\section{Особенности реализации}

\subsection{Общий поток выполнения}

Проект собирается как набор статических библиотек через CMake. Точка входа~---
единый исполняемый файл \texttt{graph\_main}, который инициализирует
\texttt{Menu} и запускает цикл ввода. Зависимости между модулями
однонаправленные: лабораторные зависят от \texttt{common}, \texttt{common}
не знает ни о каких лабораторных.

Типичный сценарий выполнения выглядит так. Пользователь выбирает пункт меню,
\texttt{Menu} вызывает соответствующий метод \texttt{Runner}'а из нужной лабы.
\texttt{Runner} читает параметры из консоли, делегирует работу алгоритмическому
классу (\texttt{Generator}, \texttt{ShimbellMethod} и т.\,д.), затем передаёт
результат в \texttt{FileHandler} для сохранения и в \texttt{Visualizer} для
отрисовки. \texttt{Visualizer} вызывает Python-скрипт через \texttt{system()},
передавая пути к файлам как аргументы командной строки. Python читает данные
из \texttt{assets/txt/} и сохраняет картинку в \texttt{assets/png/}.

\begin{figure}[h]
\centering
\begin{verbatim}
main.cc -> Menu -> Runner (lab1/lab2/...) -> Generator ->
    -> Algorithms -> FileHandler -> assets/txt/ -> 
        -> Visualizer -> Python -> assets/png/
\end{verbatim}

\caption{Поток выполнения}
\end{figure}

\subsection{Модуль common}

Модуль содержит весь код, который переиспользуется несколькими лабораторными:
базовый граф, генератор, визуализатор, файловый обработчик, утилиты и меню.
Сборка через CMake в статическую библиотеку \texttt{graph\_common}.

\subsubsection{GraphBase: шаблонный граф}

\texttt{GraphBase<VertexT, EdgeT>}~--- базовый шаблонный класс, параметризованный
типами вершины и ребра.

\begin{lstlisting}[style=cppstyle, caption={Объявление GraphBase (GraphBase.h)}]
template <VertexType VertexT, EdgeType EdgeT>
class GraphBase {
    unordered_map<uint64_t, unique_ptr<VertexT>> m_vertices_;
    unordered_map<uint64_t, EdgeT> m_edges_;
    bool is_directed_ = false;
    ...
};
\end{lstlisting}

Класс служит основой для \texttt{Graph} (\texttt{GraphBase<Vertex, EdgeData>})
и \\ \texttt{FlowNetwork} (\texttt{GraphBase<FlowVertex, FlowEdge>}).

\paragraph{Хранение данных.} Вершины и рёбра лежат в хэш-таблицах с ключом
$\mathit{key}(from, to) = (from \ll 32)\;|\;to$. Для неориентированных
графов пара нормализуется перед хешированием: $from \leq to$. Это даёт
$O(1)$ проверку наличия ребра и автоматически обеспечивает симметрию.

\paragraph{Список инцидентности.} Смежные вершины хранятся в каждом объекте
вершины в \texttt{vector<int> m\_neighbors\_}. При вызове \texttt{addEdge}
обновляются и хэш-таблица рёбер, и список соседей. Запрос
\texttt{neighbors(v)} возвращает этот вектор за $O(\deg v)$. Такая
двойная структура~--- хэш по рёбрам плюс список соседей в вершине~---
позволяет не выбирать между скоростью обхода и скоростью проверки.

\paragraph{RAII и владение.} RAII (\textit{Resource Acquisition Is Initialization})~---
идиома C++, при которой ресурс захватывается в конструкторе и освобождается
в деструкторе, время жизни ресурса привязано к времени жизни объекта.
Вершины хранятся как \texttt{unique\_ptr<VertexT>}~---
владение однозначно, удаление происходит автоматически при разрушении
контейнера, утечки памяти исключены на уровне типа.

Ребро хранится как значение \texttt{EdgeT}~--- копирование дёшево
(структура из двух \texttt{int} и \texttt{double}). Вершина~--- указатель, потому что
\texttt{VertexT} содержит \texttt{vector<int> m\_neighbors\_},
который растёт динамически; хранение по значению в \texttt{unordered\_map}
привело бы к инвалидации ссылок при рехешировании.

Все немодифицирующие методы помечены \texttt{[[nodiscard]]}~---
атрибут стандарта C++17, при котором компилятор предупреждает об игнорировании
возвращаемого значения:

\begin{lstlisting}[style=cppstyle, caption={Публичный интерфейс GraphBase (GraphBase.h)}]
[[nodiscard]] bool hasVertex(int id) const noexcept;
[[nodiscard]] bool hasEdge(int from, int to)  const;
[[nodiscard]] std::optional<EdgeT> getEdge(int from, int to);
[[nodiscard]] std::vector<int> vertexIds() const;
[[nodiscard]] std::vector<EdgeT> edges() const;
[[nodiscard]] std::vector<int> neighbors(int id) const;
[[nodiscard]] int degree(int v) const;
\end{lstlisting}

Возвращаемые значения с неопределённым результатом (отсутствующее ребро)
возвращаются через \texttt{std::optional<EdgeT>}~--- явный способ
выразить «может не быть», без исключений и без неявного кода ошибки.

\subsubsection{Graph (DAG)}

\texttt{Graph} наследует \texttt{GraphBase<Vertex, EdgeData>}.

\begin{lstlisting}[style=cppstyle, caption={Graph, Vertex, EdgeData (Graph.h)}]
struct EdgeData {
    int from{-1};
    int to{-1};
    double weight{1.0};

    EdgeData() = default;
    EdgeData(int f, int t, double w) : from(f), to(t), weight(w) {}
    bool operator<(EdgeData const& other) const { return weight < other.weight; }
};

class Vertex {
public:
    explicit Vertex(int id) : m_id_(id) {}
    [[nodiscard]] int id() const noexcept { return m_id_; }
    void addNeighbor(int neighborId);
    [[nodiscard]] std::vector<int> neighbors() const;
private:
    int m_id_;
    std::vector<int> m_neighbors_;
};

class Graph : public GraphBase<Vertex, EdgeData> {
public:
    explicit Graph(bool isDirected = false) : GraphBase(isDirected) {}
    bool addEdge(int from, int to, double weight);
    [[nodiscard]] std::optional<double> getEdgeWeight(int from, int to) const;
    [[nodiscard]] std::vector<std::pair<int, double>> neighbors(int id) const;
    [[nodiscard]] std::vector<int> getNeighbors(int id) const;
    void printGraphInfo() const;
};
\end{lstlisting}

\texttt{EdgeData} хранит поля \texttt{from}, \texttt{to} и \texttt{weight}.
Оператор \texttt{<} по весу позволяет сортировать рёбра без дополнительных
компараторов~--- используется в алгоритмах, требующих упорядочения рёбер.

\texttt{Vertex} хранит идентификатор и список смежных вершин
\texttt{m\_neighbors\_}. Метод \texttt{addNeighbor} вызывается из
\texttt{Graph::addEdge}: для неориентированного графа смежность
добавляется в обоих направлениях~--- ребро регистрируется в хэш-таблице
через \texttt{GraphBase::addEdge}, а входящая вершина получает обратную
запись через явный вызов \texttt{m\_vertices\_[to]->addNeighbor(from)}.

\texttt{Graph} расширяет интерфейс базового класса тремя методами. \\
\texttt{getEdgeWeight(from, to)} возвращает \texttt{optional<double>}~---
вес ребра или \texttt{nullopt} если ребра нет.
\texttt{neighbors(id)} возвращает список смежных вершин с весами
\texttt{vector<pair<int,double>>}, объединяя данные из списка соседей
и хэш-таблицы рёбер за один проход.
\texttt{getNeighbors(id)} делегирует в \texttt{GraphBase::neighbors}
и возвращает только идентификаторы~--- используется там, где вес не нужен.

\subsubsection{Generator: генерация графов}

Генератор использует шаблонный приватный метод,
принимающий лямбды \texttt{addEdge} и \texttt{genWeight}. Это позволяет
переиспользовать один алгоритм структурной генерации для разных типов
графов и способов генерации весов. Публичные методы лишь собирают
нужные лямбды и передают их в шаблонный метод.

Для управления режимами используются \texttt{enum class}~--- исключают
неявное преобразование в \texttt{int}:

\begin{lstlisting}[style=cppstyle, caption={Перечисления параметров генератора (Generator.h)}]
enum class EdgeCountDist { Uniform, TruncatedNormal };
enum class WeightSign    { Positive, Negative, Mixed };
\end{lstlisting}

\paragraph{Степенная генерация: generateRiceGraphByDegrees.}
Метод строит граф в три шага, связанных через \texttt{buildByRice}.

\textbf{Шаг~1: число рёбер (\texttt{sampleEdgeCount}).}

\textbf{Вход:} $m_{\min} = n-1$, $m_{\max} = \frac{n(n-1)}{2}$,
режим \texttt{EdgeCountDist}.

\textbf{Выход:} \texttt{int}~--- число рёбер $m$.

\textbf{Описание:} \texttt{Uniform}~--- равномерная выборка из
$[m_{\min}, m_{\max}]$. \texttt{TruncatedNormal}~--- нормальное
распределение $\mathcal{N}(\mu, \sigma^2)$ с $\mu = (m_{\min} + m_{\max})/2$
и $\sigma = (m_{\max} - m_{\min})/6$, усечённое методом отклонения.

\begin{lstlisting}[style=cppstyle, caption={Выбор числа рёбер (Generator.cc)}]
int Generator::sampleEdgeCount(int minE, 
                               int maxE, 
                               EdgeCountDist dist) 
{
    if (minE >= maxE) return minE;
    if (dist == EdgeCountDist::Uniform)
        return randomInt(minE, maxE);
    double mu    = (minE + maxE) / 2.0;
    double sigma = (maxE - minE) / 6.0;
    double val;
    
    do { val = mu + sigma * normal_dist_(rng_); }
    while (val < minE || val > maxE);
    
    return static_cast<int>(std::round(val));
}
\end{lstlisting}

\textbf{Шаг~2: целевые степени (\texttt{computeDegreesFromRice}).}

\textbf{Вход:} $n$, $a$, $h$, $m$.

\textbf{Выход:} \texttt{vector<int>} длины $n$.

\textbf{Описание:} Для каждой вершины сэмплируется
$x_i = \sqrt{(h + au_1)^2 + (au_2)^2}$, $u_1, u_2 \sim \mathcal{N}(0,1)$.
Вектор нормируется через softmax со стабилизацией. Степень вершины~$i$:
$q_i = p_i \cdot \mathit{cap}_i$, где $\mathit{cap}_i = n - 1 - i$~---
максимальная степень в топологическом порядке. Квоты масштабируются к
сумме $m$, целые части берутся через \texttt{floor}, остаток
распределяется жадно по убыванию дробных частей.

\begin{lstlisting}[style=cppstyle, caption={Вычисление целевых степеней (Generator.cc)}]
vector<int> Generator::computeDegreesFromRice(int n, 
                                              bool isDirected, 
                                              double a, 
                                              double h, 
                                              int totalEdges)
{
    std::vector<double> prob(n);
    for (auto& v : prob) v = riceWeight(a, h);

    double maxVal = *std::max_element(prob.begin(), prob.end());
    double sumExp = 0.0;
    for (auto& v : prob) { v = std::exp(v - maxVal); sumExp += v; }
    for (auto& v : prob) v /= sumExp;

    std::vector<int> cap(n);
    for (int i = 0; i < n; ++i) cap[i] = n - 1 - i;

    std::vector<double> quota(n);
    double sumW = 0.0;
    for (int i = 0; i < n; ++i) { quota[i] = prob[i] * cap[i]; sumW += quota[i]; }
    if (sumW > 0.0)
        for (auto& q : quota) q *= static_cast<double>(totalEdges) / sumW;

    std::vector<int> degrees(n);
    int allocated = 0;
    for (int i = 0; i < n; ++i) {
        degrees[i] = std::min((int)std::floor(quota[i]), cap[i]);
        allocated += degrees[i];
    }
    int remainder = totalEdges - allocated;
    std::vector<int> order(n);
    std::iota(order.begin(), order.end(), 0);
    std::sort(order.begin(), order.end(), [&](int a, int b) {
        return (quota[a] - std::floor(quota[a]))
             > (quota[b] - std::floor(quota[b]));
    });
    for (int idx : order) {
        if (remainder <= 0) break;
        if (degrees[idx] < cap[idx]) { ++degrees[idx]; --remainder; }
    }
    return degrees;
}
\end{lstlisting}

\textbf{Шаг~3: построение графа (\texttt{generateByDegreesTemplate}).}

\textbf{Вход:} $n$, вектор \texttt{degrees}, лямбды \texttt{addEdge}
и \texttt{genWeight}, флаг ориентированности.

\textbf{Выход:} \texttt{unique\_ptr<GraphT>}.

\textbf{Описание:} Первый проход строит остовное дерево: для вершины $i$
предок выбирается из вершин с ненулевым бюджетом степени; если таких нет~---
случайно. Второй проход добирает рёбра согласно оставшемуся бюджету.
Ациклические свойства гарантируются тем, что рёбра в обоих проходах
добавляются только от $u$ к $v$ при $u < v$.

\begin{lstlisting}[style=cppstyle, caption={Построение графа по целевым степеням (Generator.tpp)}]
template <typename GraphT, 
          typename AddEdgeFunc, 
          typename WeightGenFunc>
std::unique_ptr<GraphT> Generator::generateByDegreesTemplate(
                                    int numVertices, 
                                    std::vector<int> degrees,
                                    AddEdgeFunc addEdge, 
                                    WeightGenFunc genWeight, 
                                    bool isDirected)
{
    auto graph = std::make_unique<GraphT>(isDirected);
    for (int i = 0; i < numVertices; ++i) graph->addVertex(i);

    for (int i = 1; i < numVertices; ++i) {
        std::vector<int> with_budget;
        for (int j = 0; j < i; ++j)
            if (degrees[j] > 0) with_budget.push_back(j);
        int from = with_budget.empty()
            ? randomInt(0, i - 1)
            : with_budget[randomInt(0, (int)with_budget.size() - 1)];
        addEdge(*graph, from, i, genWeight());
        if (degrees[from] > 0) --degrees[from];
    }

    for (int i = 0; i < numVertices - 1; ++i) {
        if (degrees[i] <= 0) continue;
        std::vector<int> available;
        for (int j = i + 1; j < numVertices; ++j)
            if (!graph->hasEdge(i, j)) available.push_back(j);
        std::shuffle(available.begin(), available.end(), rng_);
        int toAdd = std::min(degrees[i], (int)available.size());
        for (int k = 0; k < toAdd; ++k)
            addEdge(*graph, i, available[k], genWeight());
    }
    return graph;
}
\end{lstlisting}

Точка входа степенной генерации:

\begin{lstlisting}[style=cppstyle, caption={generateRiceGraphByDegrees (Generator.cc)}]
std::unique_ptr<Graph> Generator::generateRiceGraphByDegrees(
                                    int numVertices, 
                                    bool isDirected, 
                                    double a, 
                                    double h,
                                    EdgeCountDist dist, 
                                    WeightSign sign)
{
    auto gen_weight = [this, a, h, sign]() {
        return applySign(riceWeight(a, h), sign);
    };
    auto add_edge = [](Graph& g, int from, int to, double w) {
        g.addEdge(from, to, w);
    };
    return buildByRice<Graph>( numVertices, isDirected, 
                               a, h, dist,
                               gen_weight, add_edge);
}
\end{lstlisting}

\subsubsection{CollectionUtils: шаблонные утилиты}

Статический класс с шаблонными методами для матриц и коллекций.

\texttt{makeMatrix<T>(rows, cols, getter)} строит матрицу любого типа,
заполняя её вызовами \texttt{getter(i, j)}. На его основе
\texttt{makeSquareMatrix} строит квадратную матрицу по списку вершин~---
используется в \texttt{FileHandler} для сохранения матриц смежности и весов.

\texttt{multiplyOptionalMatrix} выполняет матричное умножение над элементами типа \texttt{optional<double>}: \texttt{nullopt} означает отсутствие пути
и не участвует в агрегации. Компаратор передаётся параметром~---
одна функция работает как для $(\min,+)$, так и для $(\max,+)$
в методе Шимбелла.

\begin{lstlisting}[style=cppstyle, caption={Обобщённое матричное умножение (CollectionUtils.h)}]
template <typename CompareFunc>
static Matrix<std::optional<double>> multiplyOptionalMatrix(
    Matrix<std::optional<double>> const& a,
    Matrix<std::optional<double>> const& b,
    CompareFunc compare)
{
    int size = static_cast<int>(a.size());
    Matrix<optional<double>> result(
        size, vector<optional<double>>(size, nullopt)
    );
    for (int i = 0; i < size; ++i)
        for (int j = 0; j < size; ++j)
            for (int k = 0; k < size; ++k)
                if (a[i][k].has_value() && b[k][j].has_value()) {
                    double d = a[i][k].value() + b[k][j].value();
                    if (!result[i][j].has_value() || 
                        compare(d, result[i][j].value()))
                        result[i][j] = d;
                }
    return result;
}
\end{lstlisting}

\subsubsection{FileHandler: сохранение данных}

Статический класс. Все методы принимают имя файла и объект для сохранения.
\texttt{saveGraph} сохраняет список рёбер в формате \texttt{u v weight} по
одному ребру на строку~--- именно этот формат читает \texttt{GraphLoader}
в Python. \texttt{savePaths} пишет число путей первой строкой, затем
каждый путь как разделённую последовательность вершин.
Приватный шаблонный метод, предназначенный для чтение матрицы \texttt{saveSquareMatrix} используется публичными методами
\texttt{saveAdjacencyMatrix}, \texttt{saveWeightMatrix},
\texttt{saveCapacityMatrix} и \texttt{saveCostMatrix}~--- все они
отличаются только геттером, который передаётся как лямбда.

\subsubsection{Visualizer: вызов Python}

Статический класс. Каждый метод формирует список строковых аргументов
и вызывает нужный скрипт через \texttt{runPythonScript}. Вызов идёт через \texttt{system()} с построенной
командной строкой.

\texttt{draw} вызывает \texttt{plot\_network.py} с флагом \texttt{--type graph}
или \texttt{--type flow}. \texttt{drawPaths}~--- \texttt{plot\_paths.py}
с дополнительным файлом путей. \texttt{drawMatrix}~--- \texttt{plot\_matrix.py}.
\texttt{drawColoredGraph}~--- \texttt{plot\_colored\_graph.py} с файлом
раскраски вершин.

\texttt{DrawDataConfig} хранит конфигурацию путей и заголовков для каждого
пункта меню в \texttt{map<int, DrawData>}. Все пути к файлам
собраны в одном месте, что упрощает добавление нового пункта меню.

\subsection{Модуль lab1}

\subsubsection{GraphMetrics: эксцентриситеты}

\textbf{Вход:} \texttt{Graph const\&} (конструктор).

\textbf{Выход:} \texttt{MetricsResult}~--- \texttt{unordered\_map<int,double>}
эксцентриситетов, \texttt{vector<int>} центра и диаметральных вершин,
\texttt{double} радиус и диаметр.

\textbf{Описание:} Граф ациклический и невзвешенный с точки зрения метрик:
эксцентриситет считается по числу рёбер. Поэтому  используется BFS~--- приватный метод \texttt{bfsSteps},
который возвращает расстояния в шагах от источника. Сложность одного
обхода $O(V + E)$, итоговая~--- $O(V(V+E))$.

После обхода \texttt{compute} вычисляет эксцентриситет каждой вершины
как максимум расстояний по всем остальным. $+\infty$ означает, что
вершина недостижима.

\begin{lstlisting}[style=cppstyle, caption={BFS и вычисление метрик (GraphMetrics.cc)}]
unordered_map<int, double> GraphMetrics::bfsSteps(int src) const 
{
    constexpr double INF = numeric_limits<double>::infinity();
    std::unordered_map<int, double> dist;
    for (int v : m_graph_.vertexIds()) dist[v] = INF;
    dist[src] = 0.0;

    std::queue<int> q;
    q.push(src);
    while (!q.empty()) {
        int u = q.front(); q.pop();
        for (auto const& [v, w] : m_graph_.neighbors(u)) {
            if (dist[v] == INF) {
                dist[v] = dist[u] + 1.0;
                q.push(v);
            }
        }
    }
    return dist;
}

MetricsResult GraphMetrics::compute() const {
    constexpr double INF = std::numeric_limits<double>::infinity();
    auto ids = m_graph_.vertexIds();
    MetricsResult result;
    double radius = INF, diameter = 0.0;

    for (int v : ids) {
        auto dist = bfsSteps(v);
        double ecc = 0.0;
        for (int u : ids) {
            if (u == v) continue;
            double d = dist.at(u);
            if (d == INF) { ecc = INF; break; }
            ecc = std::max(ecc, d);
        }
        result.eccentricities[v] = ecc;
        if (ecc != INF) {
            radius   = std::min(radius, ecc);
            diameter = std::max(diameter, ecc);
        }
    }
    result.radius   = radius == INF ? 0.0 : radius;
    result.diameter = diameter;
    for (int v : ids) {
        double ecc = result.eccentricities[v];
        if (ecc == result.radius)   result.center.push_back(v);
        if (ecc == result.diameter) result.diametralVerts.push_back(v);
    }
    std::sort(result.center.begin(), result.center.end());
    std::sort(result.diametralVerts.begin(), result.diametralVerts.end());
    return result;
}
\end{lstlisting}

\subsubsection{ShimbellMethod: матричная степень}

\textbf{Вход:} \texttt{Graph const\&} (конструктор), \texttt{int pathLength}.

\textbf{Выход:} \texttt{ShimbellResult} --- два поля имеют тип вида
\texttt{DistanceMatrix}~$=$~\\\texttt{vector<vector<optional<double>>>}:
\texttt{min\_distances} и \texttt{max\_distances}, и \texttt{int path\_length}.

\textbf{Описание:} Начальная матрица строится через вспомогательный метод\\
\texttt{CollectionUtils::makeSquareMatrix} с геттером-лямбдой:
диагональ~--- нули, ребро~--- его вес из \texttt{getEdgeWeight},
иначе~--- \texttt{nullopt}. Матрица итеративно умножается $k-1$ раз.
Умножение вынесено в приватные методы \texttt{multiplyMin} и
\texttt{multiplyMax}, которые делегируют в
\texttt{CollectionUtils::multiplyOptionalMatrix} с соответствующим
компаратором. Использование \texttt{nullopt} 
исключает ложные пути через несуществующие рёбра. Краевой случай $k=0$ возвращает единичную матрицу. 

\begin{lstlisting}[style=cppstyle, caption={Матрица смежности и метод Шимбелла (ShimbellMethod.cc)}]
DistanceMatrix ShimbellMethod::createAdjacencyMatrix() const {
    return CollectionUtils::makeSquareMatrix<optional<double>>(
        m_vertex_ids_,
        [&](int u, int v) -> std::optional<double> {
            if (u == v) return 0.0;
            auto w = m_graph_.getEdgeWeight(u, v);
            return w.has_value() ? w : std::nullopt;
        }
    );
}

ShimbellResult const& ShimbellMethod::compute(int pathLength) {
    if (pathLength == 0) {
        DistanceMatrix identity(m_size_,
            vector<optional<double>>(m_size_, nullopt)
        );
        for (int i = 0; i < m_size_; ++i) identity[i][i] = 1;
        result_ = ShimbellResult{identity, identity, 0};
        return result_;
    }

    DistanceMatrix current_min = createAdjacencyMatrix();
    DistanceMatrix current_max = createAdjacencyMatrix();
    
    DistanceMatrix base_min    = current_min;
    DistanceMatrix base_max    = current_max;

    for (int step = 2; step <= pathLength; ++step) {
        current_min = multiplyMin(current_min, base_min);
        current_max = multiplyMax(current_max, base_max);
    }
    result_ = ShimbellResult{
        .min_distances = current_min,
        .max_distances = current_max,
        .path_length   = pathLength
    };
    return result_;
}

DistanceMatrix ShimbellMethod::multiplyMin(DistanceMatrix const& a, DistanceMatrix const& b) const {
    return CollectionUtils::multiplyOptionalMatrix(
    a, b, std::less<>()
    );
}

DistanceMatrix ShimbellMethod::multiplyMax(DistanceMatrix const& a, DistanceMatrix const& b) const {
    return CollectionUtils::multiplyOptionalMatrix(
    a, b, std::greater<>()
    );
}
\end{lstlisting}

\subsubsection{PathCounter: подсчёт маршрутов}

\textbf{Вход:} \texttt{Graph const\&} (конструктор), \texttt{int from},
\texttt{int to}.

\textbf{Выход:} \texttt{PathMatrix}~$=$~\texttt{vector<vector<int>>} --- все пути от \texttt{from} до \texttt{to}.

\textbf{Описание:} Публичный метод \texttt{getAllPaths} сначала проверяет
наличие обеих вершин через \texttt{hasVertex} и возвращает пустой результат
при отсутствии, не запуская обход. Затем вызывается рекурсивный
\texttt{backtrackAllPaths}~--- DFS с маской посещённых вершин
\texttt{unordered\_map<int,bool>}. При достижении цели текущий путь
копируется в результат; при выходе из рекурсии вершина снимается с
отметки (\emph{backtracking}). Сложность: $O(n!)$ в худшем случае.

\begin{lstlisting}[style=cppstyle, caption={Backtracking-поиск маршрутов (PathCounter.cc)}]
PathMatrix PathCounter::getAllPaths(int from, int to) {
    if (!m_graph_.hasVertex(from) || !m_graph_.hasVertex(to))
        return {};
    PathMatrix all_paths;
    std::vector<int> current_path;
    std::unordered_map<int, bool> visited;
    backtrackAllPaths(from, to, visited, 
                      current_path, all_paths
    );
    return all_paths;
}

void PathCounter::backtrackAllPaths(
    int current, int target,
    std::unordered_map<int, bool>& visited,
    std::vector<int>& currentPath, PathMatrix& allPaths)
{
    currentPath.push_back(current);
    visited[current] = true;
    if (current == target) {
        allPaths.push_back(currentPath);
    } else {
        for (auto const& [neighborId, _] : m_graph_.neighbors(current))
            if (!visited[neighborId])
                backtrackAllPaths(neighborId, target, visited, 
                                  currentPath, allPaths);
    }
    currentPath.pop_back();
    visited[current] = false;
}
\end{lstlisting}

\subsection{Модуль lab2}

\subsubsection{BFS: обход рёбер}

\textbf{Вход:} \texttt{Graph const\&} (конструктор), \texttt{int start}.

\textbf{Выход:} \texttt{BFSResult}~--- \texttt{visited\_order}
(\texttt{vector<int>}), \texttt{edges\_traversed}
(\texttt{vector<pair<int,int>>}), \texttt{parent}
(\texttt{unordered\_map<int,int>}), \\ \texttt{dist\_steps}
(\texttt{unordered\_map<int,int>}), \texttt{iterations} (\texttt{size\_t}).

\textbf{Описание:} BFS с очередью. Счётчик
\texttt{iterations} инкрементируется при каждом рассмотрении ребра $(u, v)$ и его извлечения из очереди. \texttt{dist\_steps[v] = -1}
означает недостижимость.

\begin{lstlisting}[style=cppstyle, caption={Обход рёбер в ширину (BFS.cc)}]
BFSResult BFS::traverse(int start) const {
    BFSResult result;
    for (int v : m_graph_.vertexIds()) {
        result.dist_steps[v] = -1;
        result.parent[v]     = -1;
    }
    result.dist_steps[start] = 0;
    result.visited_order.push_back(start);

    std::queue<int> q;
    q.push(start);
    while (!q.empty()) {
        int u = q.front(); q.pop();
        ++result.iterations;
        for (auto const& [v, w] : m_graph_.neighbors(u)) {
            ++result.iterations;
            if (result.dist_steps[v] == -1) {
                result.dist_steps[v] = result.dist_steps[u] + 1;
                result.parent[v]     = u;
                result.edges_traversed.emplace_back(u, v);
                result.visited_order.push_back(v);
                q.push(v);
            }
        }
    }
    return result;
}
\end{lstlisting}

\subsubsection{Dijkstra: кратчайшие пути}

\textbf{Вход:} \texttt{Graph const\&} (конструктор), \texttt{int start}.

\textbf{Выход:} \texttt{DijkstraResult}~--- \texttt{dist}
(\texttt{unordered\_map<int,double>}), \texttt{parent}
(\texttt{unordered\_map<int,int>}), \texttt{iterations} (\texttt{size\_t}).

\textbf{Описание:} Min-heap на \texttt{priority\_queue} с
\texttt{std::greater<>}. Счётчик \texttt{iterations} инкрементируется
при каждой попытке релаксации ребра и работой с кучей. Устаревшие записи в
куче отфильтровываются проверкой \texttt{d > dist[u]}.
Путь восстанавливается через \texttt{PathUtils::reconstructPath}
по массиву \texttt{parent}. Сложность: $O((V+E)\log{V})$.

\begin{lstlisting}[style=cppstyle, caption={Алгоритм Дейкстры (Dijkstra.cc)}]
DijkstraResult Dijkstra::compute(int start) const {
    DijkstraResult result;
    for (int v : m_graph_.vertexIds()) {
        result.dist[v]   = INF;
        result.parent[v] = -1;
    }
    result.dist[start] = 0.0;

    using Entry = std::pair<double, int>;
    std::priority_queue<Entry, std::vector<Entry>,
                        std::greater<Entry>> pq;
    pq.push({0.0, start});

    while (!pq.empty()) {
        auto [d, u] = pq.top(); pq.pop();
        ++result.iterations;
        if (d > result.dist[u]) continue;
        for (auto const& [v, w] : m_graph_.neighbors(u)) {
            ++result.iterations;
            double nd = result.dist[u] + w;
            if (nd < result.dist[v]) {
                result.dist[v]   = nd;
                result.parent[v] = u;
                pq.push({nd, v});
            }
        }
    }
    return result;
}
\end{lstlisting}

\subsubsection{Comparator: сравнение алгоритмов}

\textbf{Вход:} \texttt{Graph const\&}, \texttt{int start}.

\textbf{Выход:} \texttt{CompareResults}~$=$~\texttt{vector<CompareResult>},
каждый элемент содержит \texttt{algorithm\_name} и \texttt{iterations}.

\textbf{Описание:} Запускает BFS и Дейкстру на одном графе с одной
стартовой вершиной и собирает счётчики итераций в единый вектор.

\begin{lstlisting}[style=cppstyle, caption={Сравнение алгоритмов (Comparator.cc)}]
Comparator::CompareResults
Comparator::compare(Graph const& graph, int start) {
    CompareResults results;
    {
        BFS bfs(graph);
        auto r = bfs.traverse(start);
        results.push_back({"BFS", r.iterations});
    }
    {
        Dijkstra dijk(graph);
        auto r = dijk.compute(start);
        results.push_back({"Dijkstra", r.iterations});
    }
    return results;
}
\end{lstlisting}

\subsection{Модуль lab3}

\subsubsection{FlowNetwork: сеть потоков}

\textbf{Вход:} флаг ориентированности (конструктор, по умолчанию
\texttt{true}).

\textbf{Выход:} \texttt{unique\_ptr<FlowNetwork>}.

\textbf{Описание:} \texttt{FlowNetwork} наследует
\texttt{GraphBase<FlowVertex, FlowEdge>}. Структура \texttt{FlowEdge}
хранит \texttt{from}, \texttt{to}, \texttt{capacity}, \texttt{cost}
и \texttt{flow}. При добавлении ребра через \texttt{addEdge}
автоматически создаётся обратное ребро с нулевой пропускной
способностью и отрицательной стоимостью~--- это стандартное
требование алгоритмов на остаточных сетях.
Остаточная пропускная способность вычисляется как
$\mathit{res}(u,v) = \mathit{capacity}(u,v) - \mathit{flow}(u,v)$.
Метод \texttt{addFlow} обновляет поток симметрично: увеличивает
прямое ребро и уменьшает обратное.

\begin{lstlisting}[style=cppstyle, caption={Добавление ребра и обратного ребра (FlowNetwork.cc)}]
bool FlowNetwork::addEdge(int from, int to,
                          double capacity, double cost)
{
    FlowEdge forward_edge(from, to, capacity, cost, 0.0);
    FlowEdge backward_edge(to, from, 0.0, -cost, 0.0);

    bool added = GraphBase::addEdge(from, to, forward_edge);
    if (added) {
        GraphBase::addEdge(to, from, backward_edge);
        m_vertices_[to]->addNeighbor(from);
    }
    return added;
}

void FlowNetwork::addFlow(int from, int to, double flow) {
    auto* edge_forward  = getEdgeMutable(from, to);
    auto* edge_backward = getEdgeMutable(to, from);
    if (edge_forward)  edge_forward->flow  += flow;
    if (edge_backward) edge_backward->flow -= flow;
}
\end{lstlisting}

\subsubsection{MaxFlow: алгоритм Форда--Фалкерсона}

\textbf{Вход:} \texttt{FlowNetwork\&} (конструктор), \texttt{int source},
\texttt{int sink}, \\ \texttt{bool enableLogging}.

\textbf{Выход:} \texttt{double}~--- значение максимального потока.
Побочный эффект~--- вектор \texttt{m\_snapshots\_} для анимации.

\textbf{Описание:} Реализован метод Форда--Фалкерсона с BFS
для поиска увеличивающего пути.
На каждой итерации BFS ищет путь в остаточной сети~---
ребро рассматривается только если
$\mathit{res}(u,v) > 0$. Пропускная способность пути~---
минимум остаточных пропускных способностей по всем рёбрам~---
вычисляется через \texttt{PathUtils::getMinPathValue}.
Поток увеличивается через \texttt{PathUtils::forEachEdgeInPath}.
При включённом логировании после каждой итерации фиксируется
снимок состояния сети \texttt{FlowSnapshot}~--- он используется
для построения GIF-анимации роста потока.
Сложность: $O(V \cdot E^2)$.

\begin{lstlisting}[style=cppstyle, caption={Алгоритм Форда--Фалкерсона (MaxFlow.cc)}]
double MaxFlow::fordFulkerson(int source, int sink,
                              bool enableLogging)
{
    double max_flow = 0.0;
    std::unordered_map<int, int> parent;
    int step = 0;

    resetFlows(m_network_);
    m_snapshots_.clear();
    if (enableLogging) captureSnapshot(step++, max_flow);

    while (bfs(source, sink, parent)) {
        auto calc_residual = [this](int u, int v) {
            return m_network_.getResidualCapacity(u, v);
        };
        double path_flow = PathUtils<double>::getMinPathValue(
            source, sink, parent, calc_residual);

        PathUtils<double>::forEachEdgeInPath(
            source, sink, parent,
            [this, path_flow](int u, int v) {
                m_network_.addFlow(u, v, path_flow);
            });

        max_flow += path_flow;
        if (enableLogging) captureSnapshot(step++, max_flow);
    }
    return max_flow;
}
\end{lstlisting}

\subsubsection{MinCostFlow: поток минимальной стоимости}

\textbf{Вход:} \texttt{FlowNetwork\&} (конструктор), \texttt{int source},
\texttt{int sink}, \texttt{double targetFlow}.

\textbf{Выход:} \texttt{Result}~--- структура с полями \texttt{flow},
\texttt{cost}, \texttt{success} и \texttt{path} последнего найденного пути.

\textbf{Описание:} Алгоритм последовательно ищет кратчайший путь
в остаточной сети по стоимости через алгоритм Беллмана--Форда
и направляет по нему максимально возможный поток, не превышающий
$\mathit{targetFlow} - \mathit{currentFlow}$. Беллман--Форд
корректно обрабатывает отрицательные стоимости обратных рёбер.
По условию задания целевой поток берётся как $\lfloor 2/3 \cdot
\mathit{maxFlow} \rfloor$. Сложность одной итерации: $O(V \cdot E)$.

\begin{lstlisting}[style=cppstyle, caption={Поток минимальной стоимости (MinCostFlow.cc)}]
MinCostFlow::Result MinCostFlow::findMinCostFlow(
    int source, int sink, int targetFlow)
{
    Result result;
    resetFlows(m_network_);
    double current_flow = 0.0, total_cost = 0.0;
    std::unordered_map<int, double> dist;
    std::unordered_map<int, int>    parent;

    while (current_flow < targetFlow
           && bellmanFord(source, sink, dist, parent))
    {
        auto calc_residual = [this](int u, int v) {
            return m_network_.getResidualCapacity(u, v);
        };
        double path_flow = std::min(
            targetFlow - current_flow,
            PathUtils<double>::getMinPathValue(
                source, sink, parent, calc_residual));

        result.path = PathUtils<double>::reconstructPath(
            source, sink, parent);

        PathUtils<double>::forEachEdgeInPath(
            source, sink, parent,
            [this, &total_cost, path_flow](int u, int v) {
                m_network_.addFlow(u, v, path_flow);
                total_cost += path_flow
                            * m_network_.getCost(u, v);
            });

        current_flow += path_flow;
    }
    result.flow    = current_flow;
    result.cost    = total_cost;
    result.success = (current_flow == targetFlow);
    return result;
}
\end{lstlisting}

\subsubsection{PathUtils: вспомогательные операции на путях}

\texttt{PathUtils<T>}~--- шаблонный статический класс для работы
с путями, заданными массивом предшественников
\texttt{unordered\_map<int,int> parent}. Используется как в
\texttt{MaxFlow}, так и в \texttt{MinCostFlow}, а также в
\texttt{Dijkstra} из лабораторной~2.

\texttt{getMinPathValue} вычисляет минимальное значение геттера
по всем рёбрам пути~--- используется для нахождения
пропускной способности пути. \texttt{forEachEdgeInPath} применяет
произвольный функтор к каждому ребру пути~--- использует

\subsection{Визуализация: Python-скрипты}

Скрипты находятся в \texttt{scripts/} и делятся на ядро (\texttt{core/})
и скрипты отрисовки (\texttt{visualization/}). Ядро содержит три компонента: \texttt{config.py} определяет все параметры отображения: цвета, размеры, колормапы. Неизменяемые
датаклассы гарантируют, что скрипт случайно не перезапишет глобальные
настройки. \texttt{GraphLoader} читает текстовые файлы в граф \textit{networkx}. \texttt{Renderer}~--- обёртка над
matplotlib/networkx с методами \texttt{setup\_plot}, \\\texttt{compute\_layout},
\texttt{draw\_nodes}, \texttt{draw\_edges}, \texttt{draw\_labels},
\texttt{finalize}. Это убирает дублирование конфигурации в каждом скрипте.

Каждый скрипт в \texttt{visualization/} принимает аргументы через
\texttt{argparse}, читает данные через \texttt{GraphLoader}, рисует через
\texttt{Renderer} и сохраняет в файл. \texttt{plot\_network.py}
обрабатывает оба типа~--- граф и сеть потоков~--- через флаг \texttt{--type}.
Толщина рёбра для графа пропорциональна весу, для потока~--- текущему потоку;
цвет ребра потока задаётся через колормап по коэффициенту загрузки
$flow/capacity$. \texttt{plot\_paths.py} получает файл с путями,
строит словарь рёбер, входящих в пути, и выделяет их цветом: каждый путь
своим цветом из \texttt{cm.rainbow}. \texttt{plot\_matrix.py} рисует
тепловую карту через seaborn с аннотациями значений.

\subsection{Веб-интерфейс}

Веб-интерфейс реализован на Streamlit (\texttt{web/}). Приложение
(\texttt{app.py}) запускает исполняемый файл C++ как дочерний процесс
через \texttt{subprocess.Popen} с открытым \texttt{stdin}. Взаимодействие
с C++ происходит через \texttt{send\_command}~--- метод просто пишет
строку в \texttt{proc.stdin}. C++ читает её как ввод пользователя в
консольном цикле и выполняет нужное действие.

Конфигурация описана в \texttt{config.py} через датаклассы.
\texttt{LabConfig} описывает лабораторную, \texttt{ActionConfig}~--- одно
действие внутри неё: номер команды (\texttt{execute\_cmd}), список
параметров ввода (\texttt{params}) и список файлов для отображения
(\texttt{images}). Класс \texttt{GraphLabApp} в \texttt{ui.py} обходит
эту конфигурацию и рендерит виджеты Streamlit динамически.

Дополнительно для лабораторной~1 строится гистограмма весов с наложенной
теоретической PDF распределения Рэлея--Райса. PDF вычисляется напрямую в
Python через \texttt{scipy.special.i0} (функция Бесселя) по той же формуле
Вадзинского, что и в C++. Это позволяет визуально сравнить реальное
распределение сгенерированных весов с теоретическим.
